\chapter{Definitions}
\label{ch:definitions}

This chapter fixes the vocabulary. Four distinct quantities are informally
called ``time'' and three are called ``memory''; separating them is most of
the work, after which the semantics in \cref{ch:semantics} is short.

\Cref{sec:glossary} defines every concept twice, informally then formally.
\Cref{sec:memory-model} assembles the state concepts into the memory model:
what a program can store, and what a read may return.
\Cref{sec:execution-model} assembles the time and \gls{operation} concepts
into the execution model: how processes make progress against that store.

The specification is \emph{prescriptive}. Where the implementation departs
from it a marker like \devmark{dev:zero-distance} points to the deviation
register in \cref{ch:deviations}.

\section{Glossary}
\label{sec:glossary}

Each entry reads informally first, then formally. The formal definitions are
normative.

\begin{notationx}[Metavariables]\label{not:vars}
$P, Q$ range over processes; $p, q, r$ over ranks; $g$ over global element
indices; $a$ over addresses; $t$ over origin times, $\delta$ over deltas,
$\eff$ over effective times. Processes and ranks are never identified.
\end{notationx}

% ===========================================================================
\subsection{Identity and structure}
\label{sec:gloss-identity}

\term{process}{Process}
A process is one participant in a distributed simulation: a running instance
of the program, in practice one MPI process. Processes are the entities that
act --- they own memory, advance along the timeline and issue operations.
Their number is fixed for a run: DSTM performs no load balancing, spawning or
migration.

\term{rank}{Rank}
A rank is not a participant but a name for one: an integer identifying a
process uniquely. Ranks, being integers, can be ordered, indexed and arranged
in a ring; a process as such cannot. Every structure in this document is
indexed by rank.

\begin{definition}[Processes and ranks]\label{def:process}
Let $\worldsize \in \Nat$, $\worldsize \ge 1$, be the number of processes,
fixed for a \gls{session}. The \emph{rank set} is
\[
  \Ranks \defeq \{0, 1, \dots, \worldsize - 1\} \subset \Nat ,
\]
totally ordered, with \emph{ring successor} $\succr{p} \defeq (p+1) \bmod
\worldsize$. The \emph{processes} are the family
$(\proc{i})_{0 \le i < \worldsize}$ of pairwise distinct entities, with
$\Procs \defeq \{\proc{i} : i \in \Ranks\}$. The map
$\rankof{\proc{i}} \defeq i$ is a bijection $\Procs \to \Ranks$.
\end{definition}

\begin{remark}[The bijection is an assumption of this model]
\label{rem:rank-process}
$\rankof{\cdot}$ is bijective only because DSTM fixes the assignment for a
whole \gls{session}: no process is created, destroyed, migrated or
renumbered. A model admitting spawning, failure, migration, or several ranks
per process would need $\rankof{\cdot}$ partial, time-dependent or
non-injective, and every rank-indexed structure below would have to be
re-attached to processes. That generality --- logical processes, virtual
agents, load balancing --- is the extension \cite{dstm2026} defers, and
\cref{def:process} is where it would be made.
\end{remark}

\term{session}{Session}
A session is one interval of DSTM activity, opened by \code{dstm\_start} and
closed by \code{dstm\_fini}. Within it the \gls{comm-graph} and every region
size are fixed; changing either means closing the session and opening
another. A process runs sessions in sequence, never two at once.

\begin{definition}[Session]\label{def:session}
A \emph{session} is a maximal interval of execution during which $\cgraph$ and
every $S_p$ are constant. All processes enter and leave it collectively.
\end{definition}

\term{region}{Region}
\alsoterm{address-space}
A region is the flat, byte-addressable memory owned by one process, sized at
session start. Each process owns exactly one, named by that process's rank:
region identifiers and rank identifiers are the same numbers.

\begin{definition}[Region and address space]\label{def:region}
The \emph{local address space} of rank $p$ is
$\Addrof{p} \defeq \{0, \dots, S_p - 1\}$ for a region size $S_p > 0$; the
\emph{region} of rank $p$ is the corresponding partition of the global address
space $\Addr$. Both are fixed by \cref{ax:pgas}.
\end{definition}

\term{allocation}{Allocation}
Data structures carve reserved extents out of a region. An allocation is one
such reservation: a contiguous run of addresses set aside for one structure
and, while it lives, off limits to every other.

Two points are easy to get wrong. An allocation is a property of
\emph{addresses}, not \glspl{cell}: it does not vary with time, and an
allocated address is allocated in every \gls{version} alike. And what is
specified is the \emph{contract} an allocator meets, never the algorithm;
which addresses come back is deliberately open.

\begin{definition}[Allocation]\label{def:allocation}
Fix a set $\Handles$ of \emph{handles}. The allocation state of rank $p$ is a
partial function
\[
  \allocof{p} \;:\; \Addrof{p} \rightharpoonup \Handles ,
\]
mapping each reserved address to the handle reserving it. Write
$\live{p} \defeq \operatorname{dom}(\allocof{p})$ for the \emph{allocated}
addresses and $\extent{p}{h} \defeq \allocof{p}^{-1}(h)$ for the \emph{extent}
of a live handle. The state is well formed when every extent is a contiguous
interval of $\Addrof{p}$.
\end{definition}

Being a function, not a predicate, is what does the work: no address carries
two handles, so \emph{distinct allocations are disjoint automatically}. That
is what derived structures rely on, and why a handle-valued map beats a
boolean.

\begin{definition}[Allocation operations]\label{def:alloc-ops}
$\mathsf{alloc}_p(n)$ yields a fresh handle $h$ and an interval
$I \subseteq \Addrof{p}$ of length $n$ disjoint from $\live{p}$, extending
$\allocof{p}$ by $I \mapsto h$; it fails if no such interval exists.
$\mathsf{free}_p(h)$ removes $\extent{p}{h}$ from the domain. Between the two,
$\extent{p}{h}$ is constant: allocations never move or resize.
\end{definition}

\begin{remark}[The policy is not specified]\label{rem:alloc-policy}
\Cref{def:alloc-ops} says which intervals are admissible, not which is
returned. First fit, best fit, a bump pointer and block splitting all satisfy
it. The choice determines which addresses a structure occupies, never what a
\gls{cell} contains, so an allocation algorithm is not a semantic object.
\end{remark}

\term{pointer}{Pointer}
A pointer names a location in some region: which region, and how far into it.
Because region identifiers are ranks, a pointer says both where the bytes live
and which process is responsible for them --- an invariant the implementation
asserts on every operation.

\begin{definition}[Pointer]\label{def:pointer}
A \emph{pointer} is a global address $\gaddr{p}{a} \in \Addr$, well formed
when $a \in \Addrof{p}$. Write $\mathit{owner}(\pi) \defeq p$ and
$\mathit{off}(\pi) \defeq a$.
\end{definition}

% ===========================================================================
\subsection{Time}
\label{sec:gloss-time}

Four quantities are informally called ``time''. Separating them is the most
useful thing this document does.

\term{logical-time}{Logical time}
Logical time is simulation time: the timeline the simulated model advances
along, unrelated to wall-clock time or MPI progress. All processes share one
timeline and sit at their own points on it --- which is the entire purpose of
the design.

\begin{definition}[Logical time]\label{def:logical-time}
$\Time \defeq \Int$, ordered as usual. Times $\eff < 0$ are
\emph{pre-history}: no \gls{version} exists there, and a read directed at one
is answered from a caller-supplied default (\cref{def:get-or-default}).
\end{definition}

\term{origin-time}{Origin time}
The origin time of an operation is where the issuing process sits on the
timeline. It is supplied by the program, not inferred: supplying it is how a
process reports its position.

\term{delta}{Delta}
The delta is a signed displacement separating when an operation happens from
when its effect is visible. Writes carry $\delta \ge 0$, reads $\delta \le 0$.
It is the model's unit of latency, in simulated time.

\term{effective-time}{Effective time}
The effective time is the \gls{version} an operation actually touches: origin
plus delta. The memory model is indexed by it; origin time only tracks clocks.

\begin{definition}[Origin time, delta, effective time]\label{def:times}
An operation carries an \emph{origin time} $t \in \Time$ and a \emph{delta}
$\delta \in \Time$; its \emph{effective time} is $\eff \defeq t + \delta$.
Writes require $\delta \ge 0$, reads $\delta \le 0$.
\end{definition}

\term{local-clock}{Local clock}
A process's local clock is the furthest along the timeline it has declared
itself to be. It moves only forward, and only when the program issues an
operation with a larger origin time. It is what the process publishes to its
neighbours.

\begin{definition}[Local clock]\label{def:local-clock}
The \emph{local clock} $\clk{p} \in \Time$ of the process $\proc{p}$ is the
supremum of the origin times $\proc{p}$ has issued so far in the current
\gls{session}, or $0$ if none. It is nondecreasing.
\devmark{dev:clock-width}
\end{definition}

\term{known-time}{Known time}
A process cannot see its neighbours' clocks; it learns of them from messages.
The known time is its current belief about a neighbour --- always a
\emph{lower} bound, never an overestimate. That one-sidedness is what makes
the visibility rules safe.

\begin{definition}[Known time]\label{def:known-time}
For $p, q \in \Ranks$, the \emph{known time} $\kn{p}{q} \in \Time$ is
$\proc{p}$'s estimate of $\clk{q}$, initially $0$, updated only by
$\kn{p}{q} \mathrel{{:}{=}} \max(\kn{p}{q}, t)$ on receipt of a message from
$\proc{q}$ stamped $t$. It is nondecreasing.
\end{definition}

\begin{invariant}[Clock soundness]\label{inv:clock-soundness}
$\kn{p}{q} \le \clk{q}$ for all $p, q \in \Ranks$, always.
\end{invariant}

\term{drift}{Drift}
Drift is the clock skew between two processes. Tolerating it without rollback
is the point of the design; bounding it is what the \gls{comm-graph} does.

\begin{definition}[Drift]\label{def:drift}
$\drift{p}{q} \defeq \clk{p} - \clk{q}$.
\end{definition}

\term{mvt}{Minimum virtual time}
The oldest version a process still holds. A memory-management quantity, not a
semantic one: it says how far back the runtime is still paying to remember.

\begin{definition}[Minimum virtual time]\label{def:mvt}
$\mvt{p} \defeq \min \{\, \eff : \Vst{p}{\eff} \ne \Absent \,\}$, and
$\clk{p}$ when rank $p$ holds no version.
\end{definition}

% ===========================================================================
\subsection{Locality}
\label{sec:gloss-locality}

\term{comm-graph}{Communication graph}
The communication graph declares the shape of a simulation: which pairs of
ranks may ever exchange \glspl{operation}, and how long, in simulated time, an
interaction across such a pair takes to become visible. Supplied once per
\gls{session} and never changed, it is the source of every optimisation DSTM
performs --- the only reason a process can conclude that another can no longer
affect it.

\begin{definition}[Communication graph]\label{def:comm-graph}
A \emph{communication graph} is a triple $\cgraph = (\Ranks, \edges, w)$ in
which $\edges \subseteq \Ranks \times \Ranks$ is the \emph{channel} relation
and $w$ assigns each channel a weight, subject to
\begin{description}[leftmargin=2.9cm,style=nextline,font=\normalfont\bfseries]
  \item[\textnormal{(G1) Irreflexivity}] $(p,p) \notin \edges$.
  \item[\textnormal{(G2) Symmetry}] $(p,q) \in \edges \iff (q,p) \in \edges$,
    and $\lag{p}{q} = \lag{q}{p}$. \devmark{dev:asymmetric}
  \item[\textnormal{(G3) Positive integer weights}]
    $w : \edges \to \Nat_{>0}$. \devmark{dev:zero-distance}
\end{description}
\end{definition}

\term{lag}{Lag}
The lag of a channel is its weight: the number of simulated timesteps an
interaction across it takes to become visible. It is a \emph{time} separation
between two neighbours, nothing more.

\begin{convention}[Lag is not a distance]\label{conv:lag-not-metric}
$w$ is deliberately \emph{not} a metric and nothing below assumes it is. It is
defined only on channels, never between arbitrary ranks; no triangle
inequality is required, so a channel may cost more than a detour around it
without making the configuration ill formed. The only axioms are (G1)--(G3),
and only (G2) and (G3) are ever used.
\end{convention}

Of the three, (G3) carries the weight. Because $w$ is integer-valued it says
every channel costs at least one timestep, which makes every cycle in
$\cgraph$ carry positive weight --- and that is what makes the dependence
graph acyclic (\cref{thm:acyclic}) and the system deadlock-free
(\cref{thm:deadlock-freedom}). A configuration admitting a zero-weight channel
has none of DSTM's guarantees \devmark{dev:zero-distance}.

\term{reachability}{Adjacency and neighbourhood}
\alsoterm{neighbourhood}
Two ranks are \emph{neighbours} when a channel joins them directly.
Non-neighbours exchange no messages and neither constrains the other's
progress. This sparsity is what makes DSTM scale: synchronisation cost depends
on neighbourhood size, not on $\worldsize$.

\begin{definition}[Neighbourhood]\label{def:neighbourhood}
\[
  \nbr{p} \defeq \{\, q \in \Ranks : (p, q) \in \edges \,\},
  \qquad
  \diam_p \defeq \max_{q \in \nbr{p}} \lag{p}{q}
\]
with $\diam_p \defeq 0$ when $\nbr{p} = \emptyset$.
\end{definition}

\term{path-weight}{Path weight}
The path weight between two ranks is the least total lag of a chain of
channels joining them. Unlike the lag itself it is defined for every pair, and
it is the quantity that bounds how far apart two processes can drift
(\cref{prop:drift-bound}).

\begin{definition}[Path weight]\label{def:path-weight}
A \emph{path} from $p$ to $q$ is a sequence $p = z_0, \dots, z_k = q$ with
$(z_{m-1}, z_m) \in \edges$ for each $m$, of \emph{weight}
$\sum_{m} \lag{z_{m-1}}{z_m}$. The \emph{path weight}
$W : \Ranks \times \Ranks \to \Nat \cup \{\infty\}$ is the least weight of
such a path, with $\pw{p}{p} = 0$ and $\pw{p}{q} = \infty$ when no path
exists.
\end{definition}

\begin{remark}[Where a metric does belong]\label{rem:path-weight-metric}
$W$ is a metric --- a least-path weight always satisfies the triangle
inequality, by construction rather than by assumption --- whereas $w$ is not.
The distinction is exactly right: $w$ is data the application supplies and
must not be over-constrained, while $W$ is derived and is used only to bound
drift and memory. Note also that $\pw{p}{q} < \infty$ does \emph{not} make $p$
and $q$ neighbours: ranks at opposite corners of a mesh have a finite path
weight, share no channel, exchange no \glspl{heartbeat} and never appear in
each other's guards.
\end{remark}

\term{lookahead}{Lookahead}
Lookahead is the guarantee the \gls{comm-graph} buys: nothing $\proc{q}$ does
now can affect $\proc{p}$'s state earlier than $\lag{p}{q}$ into the future.
It is the same quantity conservative PDES calls lookahead \cite{chandy1979},
playing the same role --- what lets a process declare state final without
asking anyone.

\term{heartbeat}{Heartbeat}
A heartbeat carries only ``my clock has advanced to $t$''. Processes broadcast
one to their neighbours whenever their clock moves. It targets no address and
carries no data; it exists so neighbours can raise their known times and make
progress.

\begin{definition}[Heartbeat]\label{def:heartbeat}
A \emph{heartbeat} from $\proc{p}$ is a message $\langle p, t\rangle$ with
$t = \clk{p}$, delivered to every $q \in \nbr{p}$, whose sole effect is
$\kn{q}{p} \mathrel{{:}{=}} \max(\kn{q}{p}, t)$.
\end{definition}

Every message, not only a heartbeat, carries its sender's stamp and raises the
receiver's known time. A heartbeat is the degenerate message carrying nothing
else.

% ===========================================================================
\subsection{State}
\label{sec:gloss-state}

\term{cell}{Cell}
A cell is the smallest addressable unit of DSTM state, indexed by
\emph{three} coordinates: a rank, a time, and an offset. Time is part of the
location's identity, not metadata attached to it. Two writes to the same
offset at different effective times do not touch the same cell.

\begin{definition}[Cell]\label{def:cell}
$\Cells \defeq \Addr \times \Time$; we write $\cell{p}{\eff}{a}$ for the cell
at global address $\gaddr{p}{a}$ and \gls{effective-time} $\eff$
(\cref{def:cell-pgas}).
\end{definition}

\term{version}{Version}
\alsoterm{publication}
A version is the whole of one region at one \gls{effective-time} --- the cells
sharing a rank and a time. It is the unit the runtime allocates, tracks and
reclaims, created on demand by the first write targeting it, and filled with
zeroes: it does \emph{not} inherit its predecessor's contents.

\begin{definition}[Version]\label{def:version}
The \emph{version} $\ver{p}{\eff}$ is the family
$\{\, \cell{p}{\eff}{a} : a \in \Addrof{p} \,\}$ with content
$\Vbuf{p}{\eff} : \Addrof{p} \to \Bytes$ and status $\Vst{p}{\eff}$.
\end{definition}

\begin{convention}[Versions are independent, not incremental]
\label{conv:no-inherit}
$\Vbuf{p}{\eff}$ is initialised to all zeroes on creation, regardless of
$\Vbuf{p}{\eff'}$ for $\eff' < \eff$. Versions are independent snapshots
indexed by time, not deltas against an evolving state.
\devmark{dev:non-inheriting}
\end{convention}

This is the most consequential modelling decision in DSTM and the most easily
misread; \cref{sec:memory-model} takes it up.

\term{version-status}{Version status}
Every version passes through the same three stages. It is first
\emph{writable}, accepting writes and refusing reads; becomes \emph{readable}
once no further write can target it, freezing its contents; becomes
\emph{discardable} once no further read can target it, and is reclaimed. The
progression is one-way.

\begin{definition}[Status lattice]\label{def:status}
$\Status \defeq \{\Absent, \Wr, \Rd, \Dc\}$, totally ordered
$\Absent < \Wr < \Rd < \Dc$. A version's status is nondecreasing in it.
\end{definition}

\term{store}{Versioned store}
The total picture: which byte sits in every cell. This is what the memory
model is about and what an implementation must reproduce.

\begin{definition}[Versioned store]\label{def:store}
A \emph{versioned store} is a total function $\Store : \Cells \to \Bytes$.
\end{definition}

% ===========================================================================
\subsection{Operations}
\label{sec:gloss-ops}

\term{operation}{Operation}
\alsoterm{put}\alsoterm{get}
DSTM exposes exactly two data operations, a write and a read, both one-sided:
the issuing process drives it, and the target executes no application code on
its behalf. Each names a target rank, address, origin time and delta.

\begin{definition}[Operations]\label{def:operation}
An \emph{operation} is
\[
  \Put{p}{q}{t}{\delta}{a}{v}
  \qquad\text{or}\qquad
  \Get{p}{q}{t}{\delta}{a}{\ell},
\]
with $p \in \Ranks$ the \emph{issuer}, $q \in \Ranks$ the \emph{target},
$a \in \Addrof{q}$ the address, $v \in \Values$ the written value and
$\ell = |v| \ge 1$. We write $\opsrc{o}, \optgt{o}, \optime{o}, \opdelta{o}$
and $\opeff{o} = \optime{o} + \opdelta{o}$.
\end{definition}

\term{footprint}{Footprint}
The set of cells an operation touches. Because cells carry a time coordinate,
footprints do too --- which is why operations at different effective times can
never conflict however much their addresses overlap. \Cref{sec:sync} builds
the dependence analysis on this.

\begin{definition}[Footprint]\label{def:footprint}
For $o$ targeting rank $q$ at address $a$ with length $\ell$ and effective
time $\eff$,
\[
  \fp{o} \defeq \{\, \cell{q}{\eff}{a'} : a \le a' < a + \ell \,\}.
\]
$\fpw{o} \defeq \fp{o}$ if $o$ is a put and $\emptyset$ otherwise;
$\fpr{o}$ dually for gets. Operations are \emph{footprint-disjoint},
$o_1 \disjoint o_2$, when $\fp{o_1} \cap \fp{o_2} = \emptyset$.
\end{definition}

\term{issue}{Issue and completion}
\alsoterm{completion}
An operation is \emph{issued} when the program calls the API and
\emph{completes} when its effect is settled. The gap differs by operation: a
put returns immediately and completes later; a get does not return until it
completes. This is why a put carries no result and a get does.

\begin{definition}[Issue, completion]\label{def:issue-complete}
The \emph{issue} of $o$ is the instant its API call is made; its
\emph{completion} the instant its effect on the store (put) or its result
(get) is determined. Puts are \emph{asynchronous}, gets \emph{synchronous}.
\end{definition}

\term{history}{History and program order}
\alsoterm{program-order}
A history records everything a program did. Each process being
single-threaded, its own operations are totally ordered; operations of
different processes are not ordered a priori.

\begin{definition}[History, program order]\label{def:history}
A \emph{history} is a set $\Hist$ of operations with, for each $p \in \Ranks$,
a total order $\po$ on $\{\, o \in \Hist : \opsrc{o} = p \,\}$. Operations
issued by distinct ranks are unordered by $\po$.
\end{definition}

% ===========================================================================
\subsection{Coordination}
\label{sec:gloss-coord}

\term{quiescence}{Quiescence}
A system is quiescent when nothing is in flight and no process will issue
more. Detecting this is harder than detecting that every process has
\emph{finished issuing}, because an operation may still be travelling after
its issuer stopped; DSTM therefore detects the two conditions in sequence.

\begin{definition}[Quiescence]\label{def:quiescence}
A configuration is \emph{quiescent} when no message is in flight, no process
holds an incomplete operation, and none will issue another this
\gls{session}.
\end{definition}

\term{barrier}{Barrier}
DSTM's only collective inside the simulation loop, expected once at the end.
It establishes quiescence by circulating a token around the rank ring in three
rounds: the first establishes that every process has stopped issuing, the
second that every process has finished serving, the third releases everybody.
Two rounds are provably insufficient (\cref{thm:barrier-three-rounds}).

\term{termination}{Termination}
Quiescence plus teardown: after the barrier every pending-operation list is
empty and every region may be destroyed safely.

% ===========================================================================
\subsection{Derived structures}
\label{sec:gloss-derived}

The concepts above are the DSTM \emph{kernel}. Applications use a distributed
data structure layered on it; those layers add addressing discipline, not new
semantics.

\term{darray}{Distributed array (darray)}
\alsoterm{partition}
A one-dimensional array of fixed-size elements cut into contiguous slices, one
per rank. It supplies the map from a global element index to a
(rank, address) pair; reads and writes then go through the kernel unchanged.

\begin{definition}[Partition]\label{def:partition}
A \emph{partition} of $m$ elements is a vector
$(m_0, \dots, m_{\worldsize-1})$ of positive integers summing to $m$, with
\emph{offsets} $b_p \defeq \sum_{q < p} m_q$. The \emph{owner} of a global
index $g \in \{0, \dots, m-1\}$ is the unique rank $p$ with
$b_p \le g < b_p + m_p$, written $\mathit{own}(g)$.
\devmark{dev:darray-bounds}
\end{definition}

\begin{definition}[Distributed array]\label{def:darray}
A \emph{darray} is $(z, (m_p)_p, (\beta_p)_p)$ with element size $z \ge 1$, a
partition $(m_p)_p$, and $\beta_p \in \Addrof{p}$ the base of an
\gls{allocation} whose extent covers $[\beta_p, \beta_p + z\,m_p)$. Index $g$
resolves to
\[
  \mathit{res}(g) \defeq \bigl(\mathit{own}(g),\ \beta_{\mathit{own}(g)}
    + (g - b_{\mathit{own}(g)}) \cdot z\bigr).
\]
\end{definition}

The $\beta_p$ are chosen independently by each process and exchanged
collectively; the model must not assume they agree.

\term{dgraph}{Distributed graph (dgraph)}
A static directed graph with versioned state on vertices and edges. Topology
is frozen at construction; only state is versioned. An edge is owned by the
owner of its \emph{source} vertex, which puts the common ``consume my inbox
locally, produce into a neighbour's inbox'' pattern on the legal side of every
ownership boundary.

\term{topology-replica}{Topology replica}
Resolving an edge $(u,v)$ to an address requires the source rank's adjacency,
so each process keeps a copy of every neighbouring rank's topology. This
replica is static, unversioned and exchanged once at construction --- a second
class of memory, and the only one outside the versioned store, which is sound
because it is immutable.

% ===========================================================================
\subsection{Idioms}
\label{sec:gloss-idioms}

Two communication patterns recur and the API is shaped around them. They look
like opposites and are the same schedule (\cref{thm:pipeline-depth}).

\term{pull-model}{Pull model}
Each process publishes into its \emph{own} region at its current time
($\delta = 0$, local); consumers reach across into that region's \emph{past}
($\delta = -\lag{p}{q}$). Data sits still; readers reach back.

\term{push-model}{Push model}
Each process writes into its \emph{neighbours'} regions at a point in their
\emph{future} ($\delta = +\lag{p}{q}$); every process reads only its own
region at its current time ($\delta = 0$, local). Data moves; readers stay
home.

\term{mailbox}{Mailbox pattern}
Index the shared structure by \emph{(destination, source)}, so cell $(d,s)$ is
written only by $\proc{s}$ and read only by $\proc{d}$. Each cell then has
exactly one writer per effective time --- precisely the write-disjointness
determinism requires (\cref{res:write-disjoint}). Both benchmark applications
in \cite{dstm2026} use it.

\section{Memory Model}
\label{sec:memory-model}

The memory model answers one question: \emph{what may a read return?}
Messages, clocks, lags and scheduling belong to the execution model and exist
only to implement what follows.

It is built in four steps. A partitioned global address space is assumed
outright; a region is a partition of it; a \gls{cell} is a global address plus
a time coordinate; the \gls{store} maps cells to bytes. Everything
characteristic of DSTM is in the third step.

% ===========================================================================
\subsection{The partitioned global address space}
\label{sec:pgas}

DSTM does not construct its address space; it assumes one, of the kind
provided by any PGAS runtime, and adds a coordinate to it. The assumption is
this.

\begin{axiom}[Partitioned global address space]\label{ax:pgas}
There is a set $\Addr$ of \emph{global addresses} such that
\begin{description}[leftmargin=2.9cm,style=nextline,font=\normalfont\bfseries]
  \item[\textnormal{(A1) Partition}] every global address is a pair
    $\gaddr{p}{a}$ of a rank $p \in \Ranks$ and a \emph{local address}
    $a \in \Addrof{p}$, where $\Addrof{p} \defeq \{0, \dots, S_p - 1\}$ for
    some $S_p > 0$:
    \[
      \Addr \defeq \{\, \gaddr{p}{a} : p \in \Ranks,\ a \in \Addrof{p} \,\} .
    \]
  \item[\textnormal{(A2) Uniqueness}] the pair is the address: $\gaddr{p}{a} =
    \gaddr{p'}{a'}$ iff $p = p'$ and $a = a'$. Local address spaces are
    therefore disjoint as partitions of $\Addr$, and $\Addr$ is their disjoint
    union.
  \item[\textnormal{(A3) Affinity}] $\gaddr{p}{a}$ has affinity to $\proc{p}$:
    the bytes reside there, and $\proc{p}$ alone stores them.
  \item[\textnormal{(A4) One-sidedness}] any process may name any global
    address, and access it without the affine process executing application
    code.
\end{description}
\end{axiom}

\begin{definition}[Region]\label{def:region-pgas}
The \emph{region} of rank $p$ is the partition
$\{\, \gaddr{p}{a} : a \in \Addrof{p} \,\}$ of $\Addr$, identified with
$\Addrof{p}$ where no confusion arises. A \gls{pointer} is a global address;
the API encodes it as a \code{(region, offset)} pair in which the region
identifier is the rank, which is (A1) in C.
\end{definition}

(A1)--(A2) are what your intuition for a PGAS already contains: $p$ selects
the partition, $a$ indexes within it, and the pair names one byte-slot
globally. (A3) is why lags are attached to ranks rather than addresses, and
(A4) is why the interface is \code{put}/\code{get} rather than message
passing. Nothing here is specific to DSTM --- UPC++ \cite{upcxx2019} and
Charm++ \cite{charmpp1993} satisfy all four.

% ===========================================================================
\subsection{Versioning: one extra coordinate}
\label{sec:store-is-a-map}

DSTM's entire departure from \cref{ax:pgas} is to index memory by one further
coordinate.

\begin{definition}[Cell]\label{def:cell-pgas}
$\Cells \defeq \Addr \times \Time$. We write $\cell{p}{\eff}{a}$ for
$(\gaddr{p}{a}, \eff)$.
\end{definition}

\begin{definition}[Memory model]\label{def:memory-model}
The state of a DSTM \gls{session} is a \gls{store}
$\Store : \Cells \to \Bytes$ (\cref{def:store}). A read of
$\cell{q}{\eff}{a}$ returns $\Store(\cell{q}{\eff}{a})$.
\end{definition}

``Versioned memory'' suggests a mutable array plus a log of its past states,
so that version $\eff$ is what the array looked like at time $\eff$. That
reading is wrong, and every subsequent confusion follows from it. Time is a
\emph{coordinate of the cell}, on the same footing as the address, not an
annotation on a location that persists across time. Three properties fix the
consequences.

\begin{definition}[Zero default]\label{def:zero-default}
$\Store$ is total. A cell that no \gls{operation} ever writes holds the byte $0$.
\end{definition}

\begin{definition}[Independence across time]\label{def:independence}
For $\eff \ne \eff'$, the cells $\cell{q}{\eff}{a}$ and $\cell{q}{\eff'}{a}$
are distinct and unrelated. Writing one does not affect the other, and the
value of one is never derived from the other. \devmark{dev:non-inheriting}
\end{definition}

\begin{definition}[Single assignment]\label{def:single-assignment}
Each cell is written at most once over the whole execution. A program in which
two distinct writes target the same cell is not a legal DSTM program
(\cref{res:write-disjoint}).
\end{definition}

Together these put DSTM memory in \emph{dynamic single assignment} form
\cite{feautrier1991} over $\Ranks \times \Time \times \Addr$: every cell is
defined at most once and never redefined. This is the discipline SSA imposes
on scalars \cite{cytron1991ssa}, lifted to a distributed setting with
simulation time as the extra index.

\begin{remark}[The consequence programs must live with]
\label{rem:republish}
Because \glspl{version} do not inherit (\cref{def:independence}), state does
not persist across timesteps by itself: a value must be rewritten at $\eff+1$
to remain visible there. There is no read-modify-write across time and no
``current value'' of an address independent of a timestamp. Treating the store
as a per-timestep mailbox --- what both benchmark applications in
\cite{dstm2026} do --- costs nothing; treating it as long-lived state means
re-publishing.
\end{remark}

\begin{remark}[Allocation is not part of the store]
\label{rem:alloc-orthogonal}
$\Store$ is total over $\Addrof{q}$ whether or not an address is allocated
(\cref{def:allocation}). \Gls{allocation} does not decide which \glspl{cell}
\emph{exist}, but which addresses a program may legally \emph{name}. Keeping
the two apart lets the memory model stay a plain total map.
\end{remark}

\begin{definition}[Addressability]\label{def:addressability}
An operation targeting process $\proc{q}$ over the address range
$[a, a + \ell)$ is \emph{addressable} when that whole range lies within a
single extent of $\allocof{q}$:
\[
  \exists\, h .\quad [a,\, a + \ell) \subseteq \extent{q}{h} .
\]
A DSTM program may only issue addressable operations. Targeting an unallocated
address, or straddling two extents, is undefined: the store still answers, but
the byte returned belongs to another structure or to no structure at all
\devmark{dev:darray-bounds}.
\end{definition}

Addressability is a well-formedness condition, not a synchronisation
restriction --- the analogue of not dereferencing a wild pointer, prior to any
question about time. It belongs in the memory model because it concerns
\emph{which memory exists to be named}, not \emph{when it may be touched}.

\begin{remark}[Absence is not observable]
Zero-fill (\cref{def:zero-default}) makes ``nobody wrote here'' and ``somebody
wrote zero'' indistinguishable; a program needing to tell them apart must
encode presence in the value. In exchange reads never fail for want of data,
and a producer may publish only the cells it changes.
\end{remark}

% ===========================================================================
\subsection{Publication: the version life cycle}
\label{sec:publication}

\Cref{def:memory-model} describes the store as though it existed at once. An
implementation must instead decide, per version, the moment after which no
further write can arrive, since only then may it answer a read. That decision
--- promotion from writable to readable --- is the pivot of the design. Status
(\cref{def:status}) moves monotonically through four values.

\begin{figure}[htbp]
\centering
\resizebox{\textwidth}{!}{%
\begin{tikzpicture}[
  node distance=26mm,
  every node/.style={font=\small},
  st/.style={draw,rounded corners=2pt,minimum width=24mm,minimum height=9mm,
             align=center,fill=resbg,draw=dstmblue,thick},
  >={Stealth[round]},
]
  \node[st,fill=white,draw=dstmgrey] (a)             {$\Absent$\\[-1pt]\scriptsize absent};
  \node[st,right=of a]               (w)             {$\Wr$\\[-1pt]\scriptsize writable};
  \node[st,right=of w]               (r)             {$\Rd$\\[-1pt]\scriptsize readable};
  \node[st,right=of r,fill=devbg,draw=dstmred] (x)   {$\Dc$\\[-1pt]\scriptsize discardable};

  \draw[->,thick] (a) -- node[above,font=\scriptsize]{first write}
                          node[below,font=\scriptsize]{\emph{create}} (w);
  \draw[->,thick] (w) -- node[above,font=\scriptsize]{no write can arrive}
                          node[below,font=\scriptsize]{\emph{publish}} (r);
  \draw[->,thick] (r) -- node[above,font=\scriptsize]{no read can arrive}
                          node[below,font=\scriptsize]{\emph{retire}} (x);
  \draw[->,thick,dashed,draw=dstmred] (x.south) .. controls +(0,-9mm) and +(0,-9mm) ..
        node[below,font=\scriptsize]{\emph{evict} (reclaim storage)} (a.south);

  \node[below=13mm of w,font=\scriptsize,align=center,text=dstmgrey]
       {writes accepted\\reads blocked};
  \node[below=13mm of r,font=\scriptsize,align=center,text=dstmgrey]
       {contents frozen\\reads served};
\end{tikzpicture}}
\caption{The version status automaton. Progression is one-way; the dashed
edge is storage reclamation, not a semantic transition --- an evicted version
is indistinguishable from one that was never created, which is sound precisely
because \emph{retire} certifies that nobody will ask for it again.}
\label{fig:status-automaton}
\end{figure}

\begin{definition}[Publication]\label{def:publication}
A version is \emph{published} when its status reaches $\Rd$. From that instant
its content function $\Vbuf{p}{\eff}$ is immutable.
\end{definition}

Both promotions carry the correctness burden, and each is a claim about the
future:

\begin{description}[leftmargin=1.4cm,style=nextline,font=\normalfont\bfseries]
  \item[\emph{publish}] asserts that \emph{no operation that will ever be
    issued can write to this version}. Getting this wrong loses an update: a
    reader observes a version that a later write was still entitled to modify.
  \item[\emph{retire}] asserts that \emph{no operation that will ever be
    issued can read this version}. Getting this wrong either deadlocks a
    reader --- the version it wants is gone and nothing will recreate it ---
    or, worse, lets a subsequent write resurrect the version and serve
    a spuriously zero-filled snapshot.
\end{description}

Neither claim can be checked directly, because a process cannot see the future
and cannot see its neighbours. Both are instead \emph{deduced} from the
\gls{comm-graph} and the \glspl{known-time}, and the deduction is sound only if the
program honours the restrictions of \cref{ch:restrictions}. The exact guards are
given in \cref{sec:l1-rules}; their soundness is
\cref{thm:write-closure,thm:discard-safety}.

\begin{remark}[Eviction is an optimisation, promotion is not]
Reclaiming a $\Dc$ version changes no observable behaviour, so an
implementation with unbounded memory could skip eviction. \emph{Retire} itself
is not optional, since it licenses eviction. This matters for
\cref{thm:space-bound}: the \gls{footprint} is bounded because retirement is
bounded, not because eviction is eager.
\end{remark}

% ===========================================================================
\subsection{What the memory model does not provide}
\label{sec:non-guarantees}

The model is unusually strong for a distributed memory: a read returns
exactly the byte written, with no coherence protocol, fences or visibility
ordering. That strength is bought by the restrictions placed on the program,
and comes with sharp limits.

\begin{description}[leftmargin=0.9cm,style=nextline,font=\normalfont\bfseries]
\item[No atomics, no read-modify-write.]
  There is no compare-and-swap, no fetch-and-add and no way to read a cell and
  write it based on what was read, at the same effective time. The phase
  discipline of \cref{sec:publication} forbids it structurally: by the time a
  version can be read it can no longer be written.

\item[No ordering between distinct cells.]
  Two writes by the same process to different cells are unordered with respect to
  each other in the store, even when one precedes the other in \gls{program-order}.
  Ordering exists only through effective time. A program needing ``$x$ before
  $y$'' must place them at different \glspl{effective-time}.

\item[No conflict detection.]
  Violating single assignment (\cref{def:single-assignment}) is undefined
  behaviour, not an error. There is no abort, no retry and no
  last-writer-wins rule to fall back on \devmark{dev:write-disjoint}.

\item[No partial-failure model.]
  \Glspl{process} do not fail. Messages are not lost, duplicated or reordered in a way
  that matters. This is an MPI-level assumption inherited wholesale.

\item[No dynamic topology.]
  Neither the rank set, the \gls{comm-graph}, the \gls{region} sizes nor a \glslink{dgraph}{dgraph}'s
  edge set may change within a session (\cref{def:session}).
\end{description}

% ===========================================================================
\subsection{Relation to neighbouring memory models}
\label{sec:model-comparison}

DSTM's store differs from familiar designs in one place: what the version
index means, and therefore what makes a version final.

\begin{table}[htbp]
\centering
\small
\begin{tabularx}{\textwidth}{@{}l >{\raggedright\arraybackslash}X
                                 >{\raggedright\arraybackslash}X
                                 >{\raggedright\arraybackslash}X@{}}
\toprule
 & \textbf{DSTM} & \textbf{MVCC} \cite{bernstein1981}
 & \textbf{Time Warp} \cite{jefferson1985} \\
\midrule
Version index
 & simulation time, chosen by the program
 & commit timestamp, chosen by the system
 & send/receive time, with saved state per event \\[2pt]
Version content
 & independent snapshot, zero-filled
 & delta against the previous committed value
 & full checkpoint of process state \\[2pt]
Write conflict
 & forbidden by contract; undefined if violated
 & detected, resolved by abort and retry
 & tolerated, resolved by rollback and antimessages \\[2pt]
Finality of a version
 & deduced from channel lags and neighbour clocks
 & established at commit
 & never; any version may be undone \\[2pt]
Reclamation
 & lag-bounded retirement
 & garbage collection of unreachable snapshots
 & fossil collection below global virtual time \\[2pt]
Read of unwritten cell
 & zero
 & last committed value
 & value in the checkpointed state \\[2pt]
Cost of drift
 & extra live versions (bounded, \cref{thm:space-bound})
 & longer snapshot chains, more aborts
 & rollback and re-execution \\
\bottomrule
\end{tabularx}
\caption{DSTM against multi-version concurrency control and optimistic
simulation. The decisive row is \emph{finality}: MVCC learns a version is
final by observing a commit that has already happened, Time Warp never learns
it at all, and DSTM \emph{predicts} it from static topology plus a lower bound
on neighbour progress.}
\label{tab:model-comparison}
\end{table}

The contrast with a partitioned global \gls{address-space} such as UPC++
\cite{upcxx2019} is of a different kind. PGAS supplies one-sided \code{put} and
\code{get} over remote memory, exactly as DSTM does, but has no version index
at all: a \code{put} overwrites, and the application is responsible for
establishing that no reader is concurrently observing the location. DSTM's
addition is precisely the time coordinate and the runtime-enforced
\gls{publication} discipline that comes with it. Everything else in the interface is
recognisably PGAS.

\section{Execution Model}
\label{sec:execution-model}

The memory model says what the store contains. The execution model says how
\glspl{process} walk along the timeline reading and writing it, and what each
must do so its neighbours can make progress.

% ===========================================================================
\subsection{Processes and the shape of a program}
\label{sec:exec-shape}

A DSTM program is one sequential program executed by every process. Each runs
a single thread of control, occupies one point on the timeline at a time, and
hosts no virtualised sub-entities: no logical processes, no migratable agents.
Agents, if wanted, are built above DSTM with their state in the \gls{store}.

The canonical shape is a loop over ticks:

\begin{lstlisting}
for (dstm_time_t t = 0; t < T; t++) {
        /* 1. consume: read what neighbours made visible at t */
        /* 2. compute: local work, arbitrary duration                */
        /* 3. produce: publish results, visible at t + delta         */
}
dstm_barrier(last_time);
\end{lstlisting}

There is no \gls{barrier} at the tick boundary. \Glspl{process} leave the loop
body at different wall-clock instants and reach tick $t+1$ at different
instants; nothing resynchronises them.

\begin{definition}[Program]\label{def:program}
The \emph{program} of rank $p$ is the sequence of \glspl{operation} $p$ issues,
ordered by $\po$ (\cref{def:history}). A program is \emph{well clocked} when
that sequence has nondecreasing \glspl{origin-time}.
\end{definition}

The runtime does not check well-clockedness \devmark{dev:monotone-origin}. It
is a program restriction (\cref{res:monotone}), and essentially every
guarantee here is conditional on it.

% ===========================================================================
\subsection{Clocks and their publication}
\label{sec:exec-clocks}

Position on the \gls{logical-time} timeline is communicated implicitly: there
is no ``advance to time $t$'' call. A process declares where it is by issuing
an operation with origin time $t$; when that exceeds the clock, the runtime
advances it and broadcasts a heartbeat to the \gls{neighbourhood}.

\begin{definition}[Clock advance]\label{def:clock-advance}
Issuing an operation with origin time $t > \clk{p}$ has the effect
\[
  \clk{p} \mathrel{{:}{=}} t,
  \qquad
  \text{heartbeat } \langle p, t \rangle \text{ sent to every } q \in \nbr{p}.
\]
Issuing with $t \le \clk{p}$ leaves the clock unchanged and sends nothing.
\end{definition}

Two things follow.

First, \textbf{a silent process blocks its neighbours.} A process that stops issuing
operations stops advancing its clock, so its neighbours' \glspl{known-time} for it
stop rising, so the \glspl{version} those neighbours hold cannot be published. A process
with nothing to say must still say it, by continuing to issue at increasing
origin times or by entering the \gls{termination} barrier.

Second, \textbf{the clock tracks origin time, not \gls{effective-time}.} A put with
origin $t$ and \gls{delta} $\delta$ advances the clock to $t$, not to $t+\delta$
\devmark{dev:clock-conflation}. This is the correct choice --- the clock's job
is to report where the \emph{issuer} is, and the issuer is at $t$ --- but it
means the largest effective time a process has produced can exceed its clock by
an unbounded amount, which is exactly what \cref{res:horizon} constrains.

\begin{definition}[Known-time update]\label{def:known-update}
Every message carries its sender's identity $q$ and a stamp $t \le \clk{q}$.
On receipt at $p$: $\kn{p}{q} \mathrel{{:}{=}} \max(\kn{p}{q}, t)$.
\end{definition}

The stamp is a \emph{lower} bound on the sender's clock, not necessarily equal
to it: a reply to a read is stamped with the read's \gls{effective-time}
rather than the responder's clock. Taking the maximum makes known times
monotone, and \cref{inv:clock-soundness} makes them sound. Underestimating is
harmless where overestimating would not be:

\begin{proposition}[Conservatism of known times]\label{prop:conservative}
Both promotion guards of \cref{sec:publication} are monotone in the known
times: lowering any $\kn{p}{q}$ can only prevent a promotion, never cause one.
Consequently an implementation that underestimates neighbour clocks is
\emph{safe but possibly slow}, whereas one that overestimates them is
\emph{unsound}.
\end{proposition}

% ===========================================================================
\subsection{Issue, completion and blocking}
\label{sec:exec-blocking}

The two \glspl{operation} differ in the gap between \gls{issue} and
\gls{completion}, and the asymmetry is deliberate.

\begin{description}[leftmargin=0.9cm,style=nextline,font=\normalfont\bfseries]
\item[A put is asynchronous.]
  It returns as soon as the payload has been copied out of the caller's
  buffer. It carries no result and cannot report a delivery failure. Its
  effect on the store is settled later, possibly long after the issuing process
  has moved on. A program therefore has no way to observe that its own put has
  landed, other than by reaching \gls{quiescence}.

\item[A get is synchronous.]
  It does not return until the requested version has been published and its
  bytes copied into the caller's buffer. While waiting it drives the runtime's
  progress engine, so a blocked process continues to serve its neighbours'
  requests, forward messages and promote its own versions. This is what keeps
  a blocked get from becoming a deadlock in the common case.
\end{description}

\begin{remark}[Blocking is where drift is absorbed]
A get is the only place a DSTM program waits, and its duration is exactly how
far the reader has run ahead of the producer it reads. Versioning, lag
bounds and \glspl{heartbeat} all exist to make that wait as short as the
application's real data dependence requires.
\end{remark}

% ===========================================================================
\subsection{The task grid and drift}
\label{sec:exec-grid}

Picture an execution as a grid: one row per rank, one column per tick, one
\emph{task} per \gls{cell}. The task $\task{p}{t}$ is the work $\proc{p}$
performs at origin time $t$ --- one loop iteration. \Glspl{operation} create
edges between tasks, weighted by their \gls{delta}.

\begin{definition}[Task]\label{def:task}
A \emph{task} $\task{p}{t}$ is the set of operations issued by $p$ with origin
time $t$, together with the local computation between them. Its \emph{work}
$\work{p}{t} \ge 0$ is the wall-clock duration of that computation.
\end{definition}

\begin{figure}[htbp]
\centering
\begin{tikzpicture}[
  x=17mm, y=-12mm,
  every node/.style={font=\scriptsize},
  tk/.style={draw,rounded corners=1.5pt,minimum width=11mm,minimum height=6mm,
             fill=resbg,draw=dstmblue,inner sep=1pt},
  >={Stealth[round]},
]
  % tick axis
  \foreach \t in {0,1,2,3,4} {
    \node[text=dstmgrey] at (\t,-0.75) {$t=\t$};
  }
  % process rows
  \foreach \r/\name/\last in {0/r_0/4, 1/r_1/3, 2/r_2/2, 3/r_3/3} {
    \node[text=dstmgrey,anchor=east] at (-0.55,\r) {$\name$};
    \draw[dstmgrey!40] (-0.35,\r) -- (4.6,\r);
    \foreach \t in {0,...,\last} { \node[tk] (n\r-\t) at (\t,\r) {$\t$}; }
  }
  % put edges: forward in time, to a neighbour  (solid)
  \draw[->,thick,dstmblue] (n0-1) -- node[above right,inner sep=1pt]{$+1$} (n1-2);
  \draw[->,thick,dstmblue] (n1-1) -- node[above right,inner sep=1pt]{$+1$} (n2-2);
  \draw[->,thick,dstmblue] (n3-1) -- node[below right,inner sep=1pt]{$+1$} (n2-2);
  \draw[->,thick,dstmblue] (n0-2) -- (n3-3);
  % get edges: backward in time, from a neighbour (dashed)
  \draw[->,thick,dashed,dstmred] (n1-2) to[bend left=18] (n0-3);
  \draw[->,thick,dashed,dstmred] (n2-2) to[bend right=18] (n3-3);
  % drift bracket
  \draw[decorate,decoration={brace,amplitude=4pt},dstmamber,thick]
       (2.55,2.35) -- (4.45,2.35)
       node[midway,below=4pt,text=dstmamber]{drift of $r_0$ over $r_2$};
\end{tikzpicture}
\caption{An execution as a task grid, for a four-process ring with all lags
equal to $1$. Solid edges are puts (forward in time, into a neighbour's
region); dashed edges are gets (backward in time, out of a neighbour's
region). Process $r_0$ has reached tick $4$ while $r_2$ is still at tick $2$:
they have drifted apart by two ticks and neither has rolled back or waited at
a barrier. The edges, not the rows, are what constrain the schedule.}
\label{fig:task-grid}
\end{figure}

\Cref{fig:task-grid} makes the central claim visible: rows are independent,
and the only coupling is the edges, every one of which spans at least one tick
because every \gls{lag} is at least one. A process is never blocked by
\emph{where} its neighbours are, only by \emph{what} they have already
produced.

Drift is not bounded by the model, but it is bounded by the
\gls{comm-graph} whenever the program communicates each tick. That hypothesis
is worth naming.

\begin{definition}[Coupled schedule]\label{def:coupled}
A schedule is \emph{coupled} when, at every tick $t$ it executes, each process
$\proc{p}$ completes a \gls{get} whose target \gls{version} cannot be
published until $\clk{q} \ge t - \lag{p}{q}$ for every $q \in \nbr{p}$. Write
$T_p$ for the greatest tick $\proc{p}$ has completed.
\end{definition}

Both standard idioms are coupled, for slightly different reasons. Under the
\gls{push-model} a process reads its own region at $\delta = 0$, so publishing
$\ver{p}{t}$ requires $\kn{p}{q} + \lag{p}{q} > t$ directly. Under the
\gls{pull-model} it reads $\ver{q}{t - \lag{p}{q}}$, whose publication
requires $\clk{q} \ge t - \lag{p}{q}$ by the local-clock clause.
\Cref{prop:drift-bound} turns this into a bound computed from $\cgraph$ alone,
and \cref{cor:space-drift} turns that into a bound on memory.

% ===========================================================================
\subsection{Session lifecycle}
\label{sec:exec-lifecycle}

\begin{figure}[htbp]
\centering
\resizebox{\textwidth}{!}{%
\begin{tikzpicture}[
  node distance=6mm and 9mm,
  every node/.style={font=\scriptsize},
  ph/.style={draw,rounded corners=2pt,minimum height=8mm,align=center,
             inner xsep=4pt,draw=dstmblue,fill=resbg},
  >={Stealth[round]},
]
  \node[ph] (init)  {\code{dstm\_init}};
  \node[ph,right=of init]  (start) {\code{dstm\_start}$(\Dmat)$};
  \node[ph,right=of start] (mk)    {\code{darray/dgraph}\\\code{\_create}};
  \node[ph,right=of mk]    (loop)  {tick loop\\\code{put}/\code{get}};
  \node[ph,right=of loop]  (fini)  {\code{dstm\_fini}$(\ell)$};
  \node[ph,right=of fini]  (final) {\code{dstm\_finalize}};
  \foreach \a/\b in {init/start,start/mk,mk/loop,loop/fini,fini/final}
    { \draw[->,thick] (\a) -- (\b); }
  \draw[->,thick,dstmteal] (fini.south) .. controls +(0,-7mm) and +(0,-7mm) ..
       node[below,text=dstmteal]{further sessions} (start.south);

  \begin{scope}[on background layer]
    \node[fit=(start)(fini),draw=dstmteal,dashed,rounded corners,
          inner sep=5pt,label={[text=dstmteal,font=\scriptsize]above:session}] {};
  \end{scope}
\end{tikzpicture}}
\caption{Session lifecycle. The \gls{comm-graph} and every region are fixed
between \code{dstm\_start} and \code{dstm\_fini}; changing topology means
closing the session and opening another.}
\label{fig:lifecycle}
\end{figure}

The restrictions attach to different scopes: \cref{res:collective} to the
whole process lifetime; \crefrange{res:monotone}{res:phase} and
\cref{res:write-disjoint} to a single \gls{session}; \cref{res:terminate} to
the transition out of one.

% ===========================================================================
\subsection{The two synchronisation idioms}
\label{sec:exec-idioms}

The API admits many \gls{delta} choices; applications use two, and both
benchmark applications in \cite{dstm2026} use one of them exclusively.

\begin{center}
\small
\begin{tabular}{@{}lll@{}}
\toprule
 & \textbf{Pull} & \textbf{Push} \\
\midrule
Producer $p$ writes & own region, $\delta = 0$
                    & neighbour $q$'s region, $\delta = +\lag{p}{q}$ \\
Consumer $q$ reads  & $p$'s region, $\delta = -\lag{p}{q}$
                    & own region, $\delta = 0$ \\
Data movement       & on demand, at read time & eager, at write time \\
Messages per pair   & two (request, reply)    & one \\
\bottomrule
\end{tabular}
\end{center}

They look like opposites, but the value produced by $\proc{p}$ at tick $t$
becomes available to $\proc{q}$ at tick $t + \lag{p}{q}$ in both cases: the
\gls{lag}, not the direction of data movement, sets the pipeline depth.
\Cref{thm:pipeline-depth} states this precisely. The choice is therefore a
communication-cost decision --- one message versus two --- never a semantic
one.
